%% ma: L12-B06-0079
%% lop: 12
%% chuong: 2
%% bai: 6
%% dang: 12.6.13
%% loai: TN4
%% muc: VD
%% dap_an: "A"
%% nguon: "CĐ2.1.pdf giáo viên gửi 27/10/2026 (M07, Lê Minh Khôi – Nguyễn Anh Khôi, Chương II Chủ đề 1 Vectơ trong không gian), Câu 40"
%% kien_thuc: "Bài 6 (bình phương vô hướng)"
%% trang_thai: chinh-thuc
\begin{ex}
Trong không gian, cho hai vectơ $\vec a$, $\vec b$ thỏa mãn $\left|2\vec a+\vec b\right|=12$ và $\left|\vec a-3\vec b\right|=3\sqrt{11}$. Tính giá trị của biểu thức $J=2|\vec a|^2+3|\vec b|^2$.
\choice{\True $J=90$}{$J=630$}{$J=70$}{$J=420$}
\loigiai{Đặt $x=|\vec a|^2$, $y=|\vec b|^2$, $t=\vec a\cdot\vec b$: $4x+4t+y=144$ và $x-6t+9y=99$.

Nhân phương trình đầu với $3$, phương trình sau với $2$ rồi cộng: $14x+21y=630\Rightarrow J=2x+3y=90$.}
\end{ex}
